\begin{table}[H]
\centering
\caption{So sánh trạng thái hệ thống trước và sau khi áp dụng biện pháp bảo mật}
\label{tab:before-after-security}
\renewcommand{\arraystretch}{1.35}
\begin{tabular}{|p{3.2cm}|p{5.4cm}|p{5.4cm}|}
\hline
\textbf{Nội dung đánh giá} & \textbf{Trước khi khắc phục} & \textbf{Sau khi khắc phục} \\
\hline
Public Storage Bucket &
Bucket \texttt{public-customer-data} được cấu hình công khai. Người dùng có thể mở trực tiếp đường dẫn object như \texttt{customers.csv} mà không cần xác thực. &
Bucket được chuyển sang chế độ riêng tư. Khi truy cập trực tiếp object, hệ thống từ chối yêu cầu hoặc yêu cầu thông tin xác thực hợp lệ. \\
\hline
IAM Misconfiguration &
Tài khoản người dùng thường \texttt{alice} bị cấp nhầm quyền quản trị và có thể truy cập trang \texttt{/admin}. &
Quyền truy cập được kiểm tra theo vai trò. Tài khoản \texttt{alice} không còn truy cập được trang quản trị. \\
\hline
Credential Leakage &
Thông tin nhạy cảm như mật khẩu cơ sở dữ liệu, access key và secret key được đặt trực tiếp trong cấu hình ở dạng rõ. &
Secret được tách khỏi mã nguồn và lưu dưới dạng đã mã hóa. Ứng dụng chỉ giải mã khi cần sử dụng trong môi trường runtime. \\
\hline
Khả năng giám sát &
Việc phát hiện truy cập bất thường còn phụ thuộc vào quan sát thủ công trong quá trình demo. &
Hệ thống có thể mở rộng bằng audit log, cảnh báo truy cập trái phép và theo dõi các chỉ số như MTTD, MTTR. \\
\hline
Tác động bảo mật &
Dữ liệu khách hàng và thông tin cấu hình có nguy cơ bị lộ nếu attacker truy cập được endpoint hoặc repository. &
Giảm nguy cơ rò rỉ dữ liệu, giới hạn quyền theo nguyên tắc least privilege và giảm thiểu tác động khi cấu hình bị lộ. \\
\hline
\end{tabular}
\end{table}
